% The viewers' HTTP interfaces: VitalStream's viewer API, version 1 (the authority is
% viewers/vitalstream/rust/src/server.rs, which the C++ and Python viewers match), the rules of
% Medorama's gateway for it (medorama/src/gateway.rs), the addresses of the ECG and PACS viewers
% (viewers/README.md) and the requests of their `serve` command (page_server.py).
\chapter{The viewer API}
\label{chap:viewer-api}

Each viewer is a web page. VitalStream's page talks to its own server through the
\emph{viewer API}, version 1, described in full in this chapter; the ECG and PACS viewers read
the files they are given in the browser, and are told what to show by their address. This
chapter is the reference of these interfaces: the requests of VitalStream's server, the rules
under which \md{} passes them on, the addresses of the three viewers, and the requests of the
small server that the ECG and PACS viewers' command \cmd{serve} runs. The viewers themselves
are described in \cref{chap:vitalstream}, \cref{chap:ecg-viewer} and \cref{chap:pacs-viewer};
their commands and options in \cref{chap:cli}.

\section{VitalStream's viewer API}
\label{sec:viewer-api-vitalstream}

The three VitalStream servers, \cmd{vitalstream-py}, \cmd{vitalstream-rs} and
\cmd{vitalstream-cpp}, implement the same API and give the same answers to the same requests;
a test of the repository (\file{viewers/vitalstream/tests/parity/viewer\_parity.py}) checks it.
A change that breaks the API gets a new version, under \file{api/v2/}.

\subsection{Conventions}
\label{sec:viewer-api-conventions}

\begin{itemize}
\item \textbf{Paths.} Every path is relative to where the viewer is served, so that the page
  also works under a path prefix (behind a reverse proxy, a cluster's portal, or \md{}'s
  \file{/viewer/}): \code{api/v1/config} is
  \code{http://localhost:8060/api/v1/config} for \cmd{vitalstream-rs} on its default port, and
  \code{https://portal.example.org/node/8060/api/v1/config} behind a portal.
\item \textbf{Times} are Unix seconds (UTC). A value that is not a number (NaN) is sent as
  \code{null}.
\item \textbf{Files.} Only files and folders inside the data roots (the setting
  \code{data\_roots}, \cref{chap:settings-reference}) are listed, read or opened; a path is
  checked after its symbolic links are resolved, so a link that leads outside the data roots is
  refused.
\item \textbf{Datasets.} A recording, once opened, is a \emph{dataset}, named by its
  \emph{key}: the first 20 hexadecimal digits of a hash of the index format's version, the
  recording's real path and its fingerprint (the sizes and times of its files), so that the key
  changes when the recording changes. The key is the same in every viewer and in \md{}. The
  index and the terms of the answers (channels, pieces, breaks, vital signs, alarms, markers)
  are described in \cref{chap:formats-reference}.
\item \textbf{Answers} are JSON. \textbf{Errors} are \code{\{"detail": "..."\}} with the status
  codes of \cref{tab:viewer-api-status}.
\end{itemize}

\begin{xltabular}{\linewidth}{@{}P{1.4cm}L@{}}
\caption{The status codes of the viewer API.}
\label{tab:viewer-api-status}\\
\toprule
{\tabhead Status} & {\tabhead Meaning} \\
\midrule
\endfirsthead
\toprule
{\tabhead Status} & {\tabhead Meaning} \\
\midrule
\endhead
\bottomrule
\endfoot
\code{400} & A bad request: a host the server does not answer for, a missing, repeated or
  malformed query parameter, an unknown channel, a bad body, a path that is not a recording the
  viewer reads, \code{t1} not after \code{t0}. Every viewer answers a malformed query parameter
  with \code{400}, never \code{422}. \\
\code{401} & Sign-in required (\cref{sec:viewer-api-signin}). \\
\code{403} & A path not allowed: outside the data roots (which lie inside the viewer's root; a
  path is checked after its symbolic links are resolved), a file of the record that is, or a
  folder that cannot be read; or a request from another site. \\
\code{404} & No such file or folder, no dataset indexed under that key, or no index being built
  under it (\code{cancel}). \\
\code{405} & A known path asked with another method (\code{\{"detail": "Method Not
  Allowed"\}}). \\
\code{413} & A request body over its limit. \\
\code{415} & A \code{POST} to the API whose body is not JSON. \\
\code{429} & A password sent before the delay after wrong ones has passed (the C++ viewer). \\
\code{422} & A folder over one of the limits on a folder of records, or a folder without
  records (\code{scan}, \code{open}); the \code{detail} says which limit and what to do
  instead. \\
\code{500} & A fault while reading a dataset, or another failure of the server while
  answering. \\
\end{xltabular}

\subsection{Checks on every request}
\label{sec:viewer-api-checks}

The server checks every request before it routes it, in this order. The Rust viewer is the
reference; the other viewers give the same statuses and the same checks, and the parity test
compares them.
\begin{enumerate}
\item \textbf{Host.} The \code{Host} header (or, without one, the request's authority) must
  name a host the server answers for, compared without its port and in any letter case:
  \code{localhost}, \code{127.0.0.1}, \code{::1} (\code{[::1]}), the \code{host} the server
  listens on (unless that is \code{0.0.0.0} or \code{::}), or a name in the setting
  \code{allowed\_hosts} (\code{"*"} allows any). Otherwise the answer is \code{400}. This
  defeats DNS rebinding, where another site's name is made to point at the viewer's address. A
  viewer behind a reverse proxy that passes the public host name on lists that name in
  \code{allowed\_hosts}; \md{}'s gateway sends \code{127.0.0.1:<port>}.
\item \textbf{Requests that change something} (any method other than \code{GET} and
  \code{HEAD}): an \code{Origin} header, when present, must be \code{http://} or
  \code{https://} followed by a host allowed as in 1 (no path), and \code{null} and any other
  origin are refused (\code{403}); \code{Sec-Fetch-Site: cross-site} is refused (\code{403});
  and a \code{POST} under \file{api/} must have \code{Content-Type: application/json}
  (parameters such as \code{; charset=utf-8} allowed; \code{415} otherwise). Another site's page
  cannot send such a request without the browser asking the server first, and the server never
  allows it.
\item \textbf{Sign-in} (\cref{sec:viewer-api-signin}): \code{401} without it.
\item \textbf{Routing:} an unknown path is \code{404}, and a known path with another method
  \code{405}.
\end{enumerate}

\textbf{Bodies.} \code{POST api/v1/open} and \code{POST api/v1/cancel} take a JSON object of at
most 64\,KiB, and \code{POST login} a form of at most 4\,KiB; a larger body is \code{413}.
\code{open} takes \code{path} (a string) and \code{reindex} (a boolean, default \code{false});
\code{cancel} takes \code{key} (a string). Any other type, such as \code{\{"reindex": "yes"\}},
is \code{400}, and so is a body that is not a JSON object (an array, a number); unknown fields
are ignored.

\textbf{Query parameters} are \code{application/x-www-form-urlencoded}. A parameter given twice
(\code{ch=0\&ch=1}) is \code{400}. Numbers are checked strictly, and anything else is
\code{400}:
\begin{itemize}
\item a time or other number (\code{t0}, \code{t1}, \code{t}) is a finite decimal number, with
  nothing around it (not even a space): \code{-inf}, \code{nan}, \code{0x10}, \code{1\_0} and
  \code{1e400} are refused;
\item an integer (\code{px}) is digits with an optional sign, within 64 bits;
\item \code{ch} is a comma-separated list of channel ids, digits only, each the id of a channel
  of the dataset. An id given more than once is answered once, in the order first given, and at
  most 64 distinct ids are taken;
\item \code{names} (\code{vitals}) is a \code{\textbar}-separated list of series names. A name
  given more than once is answered once, at most 512 distinct names are taken, and names the
  dataset does not have are left out of the answer;
\item in \code{window}, \code{vitals} and \code{markers}, \code{t1} must be greater than
  \code{t0}. Times far outside the recording are allowed and simply find no data.
\end{itemize}

\textbf{Answers} carry \code{X-Content-Type-Options: nosniff}. The page and the sign-in page
carry the Content Security Policy \code{default-src 'self'; script-src 'self'; style-src 'self';
img-src 'self' data:; object-src 'none'; base-uri 'none'; frame-ancestors 'self'; form-action
'self'}; they have no inline script or style (the sign-in page links \file{static/login.css}).

\subsection{Signing in}
\label{sec:viewer-api-signin}

When a password is set (with \opt{-{}-set-password} of any of the three viewers: one password,
in one file, for all of them, \cref{chap:cli}), every request except the page itself
(\code{GET /}), \code{POST login}, \code{POST logout} and \code{GET static/...} needs the
sign-in cookie, named \code{vitalstream\_<port>} after the server's port; without it, the API
answers \code{401} \code{\{"detail": "sign in required"\}}, and the page shows a sign-in form
(with the status \code{401}). The cookie is
\code{vitalstream\_<port>=<expiry>.<HMAC-SHA256(secret, expiry)>}, signed with the password
file's secret and valid for 30 days (\code{HttpOnly}, \code{Path=/}, \code{SameSite=Lax}, and
\code{Secure} when the setting \code{secure\_cookie} is \code{true}, for a viewer behind an HTTPS
proxy). A cookie whose signature is in the password file's \code{revoked} list is refused.
Without a password set, no request needs a cookie.

Passwords are checked one at a time (scrypt uses 16\,MiB). After wrong passwords from one client
address, the next password from it is checked only after a delay: $2^n-1$ seconds after $n$
wrong ones (1, 3, 7, 15\,s), at most 30\,s; the count is kept 15 minutes after the last wrong
password, and a right password clears it. The Rust and Python viewers hold such a request until
the delay has passed; the C++ viewer does not hold it: until the delay has passed it answers
\code{429} with \code{Retry-After} (the seconds left) and the sign-in page.

\apientry{GET /}{The page}
\begin{lstlisting}[style=signature]
GET /
\end{lstlisting}
The viewer's page; before sign-in (with a password set), the sign-in form, with the status
\code{401}.

\apientry{GET static/...}{The page's files}
\begin{lstlisting}[style=signature]
GET static/{path}
\end{lstlisting}
The page's scripts, styles and pictures. Need no sign-in.
\raises{\code{404} no such file.}

\apientry{POST login}{Sign in}
\begin{lstlisting}[style=signature]
POST login   (form: password=PASSWORD)
\end{lstlisting}
Checks the password sent by the sign-in form (\code{application/x-www-form-urlencoded}, at
most 4\,KiB), sets
the sign-in cookie and redirects (\code{303}) to \code{./}. A wrong password is answered, after
a pause of one second, with the form again and the status \code{401}. Without a password set,
only redirects to \code{./}.

\apientry{POST logout}{Sign out}
\begin{lstlisting}[style=signature]
POST logout
\end{lstlisting}
Removes the sign-in cookie and redirects (\code{303}) to \code{./}. When the request carries a
valid sign-in cookie, that sign-in also ends on the server: its signature is added to
\code{revoked} in the password file (\code{\{"<signature>": <expiry>\}}, kept until the
expiry), so that a copy of the cookie no longer works, also after a restart.

\subsection{The requests}
\label{sec:viewer-api-requests}

\apientry{GET api/v1/config}{The server's settings for the page}
\begin{lstlisting}[style=signature]
GET api/v1/config
\end{lstlisting}
What the page needs to know about the server.
\returns{\code{timezone} (the time zone in which times are shown), \code{data\_roots},
\code{home} (the home folder of the account that runs the viewer), \code{start} (the folder the
file browser opens at: the home folder if it lies inside the data roots, else the first data
root), \code{places} (the shortcuts of the Open dialog: \code{\{name, path\}} for the home
folder, the repository's \file{sample\_data} folder when the viewer runs from the repository,
and each data root other than the home folder and the root of the file system), \code{sep} (the
path separator, \code{/} or \code{\textbackslash}), \code{version} (of the viewer),
\code{qc\_verbose} (the setting) and \code{limits} (\code{max\_records}, \code{max\_depth},
\code{max\_scan\_s}: the limits on a folder of records, \cref{chap:recordings}).}
\begin{lstlisting}[language=JSON]
{"timezone": "America/New_York", "data_roots": ["/data/icu"], "home": "/home/reader1",
 "start": "/data/icu", "places": [{"name": "/data/icu", "path": "/data/icu"}], "sep": "/",
 "version": "0.1.0", "qc_verbose": false,
 "limits": {"max_records": 5000, "max_depth": 20, "max_scan_s": 10.0}}
\end{lstlisting}

\apientry{GET api/v1/browse}{A folder's entries}
\begin{lstlisting}[style=signature]
GET api/v1/browse?path=FOLDER&hidden=BOOL
\end{lstlisting}
The entries of a folder, for the file browser.
\begin{params}
\item[path] The folder; default the \code{start} folder of \code{config}.
\item[hidden] \code{true} (or \code{1}, \code{yes}, \code{on}): also the hidden entries (names
  that begin with \code{.}) and the segments and layout headers of the multi-segment records of
  the folder. Default \code{false}.
\end{params}
Each entry has a \code{name}, a \code{path}, a \code{type} and \code{link} (whether it is a
symbolic link):
\begin{itemize}
\item \code{dir}: a folder;
\item \code{wfdb}: a folder that holds WFDB records itself, or a WFDB collection
  (\cref{chap:recordings});
\item \code{record}: a WFDB record, its \file{.hea} file: \code{name} without \file{.hea},
  \code{kind} \code{wfdb}, \code{label} \code{WFDB record} or \code{WFDB multi-segment record};
\item \code{file}: a recording file, with its \code{kind}, the format's \code{label} and its
  \code{size} in bytes. It is an EDF or BDF file (\code{kind} \code{edf}, \code{label}
  \code{EDF}, \code{EDF+}, \code{BDF} or \code{BDF+}; \cref{sec:vitalstream-edf}), or a file
  that an installed adapter reads (\cref{chap:addons});
\item \code{broken}: a symbolic link that leads nowhere.
\end{itemize}
Not listed: the entries that lead outside the data roots, the signal files and every other file
that no reader reads, and (without \code{hidden}) the hidden entries and the segments of the
multi-segment records. Folders come first, then records and files, each by name (ignoring
case); at most 5000 entries are listed.
\returns{\code{\{path, type: "dir", kind, entries, truncated, parent\}}: \code{kind} the
folder's own kind (\code{wfdb}, or \code{null}), \code{truncated} \code{true} when entries were
left out after 5000, \code{parent} the folder above (\code{null} at the top of the data roots).
A path to a recording file (a WFDB record's \file{.hea}, an EDF or BDF file, or a file an
adapter reads) answers
\code{\{"path": PATH, "type": "file"\}}.}
\raises{\code{400} neither a folder nor a recording; \code{403} outside the data roots, or a
folder that cannot be read; \code{404} no such folder.}
\begin{lstlisting}[language=JSON]
{"path": "/data/icu/demo", "type": "dir", "kind": "wfdb", "truncated": false,
 "parent": "/data/icu",
 "entries": [{"name": "bed01", "path": "/data/icu/demo/bed01", "type": "dir", "link": false},
             {"name": "rec0001", "path": "/data/icu/demo/rec0001.hea", "type": "record",
              "kind": "wfdb", "label": "WFDB record", "link": false}]}
\end{lstlisting}

\apientry{GET api/v1/complete}{Complete a path}
\begin{lstlisting}[style=signature]
GET api/v1/complete?prefix=TEXT
\end{lstlisting}
The paths that begin with a prefix, for the path field of the Open dialog: the entries of the
prefix's folder whose names begin with its last part (ignoring case), as \code{browse} lists
them; hidden names only when the last part begins with \code{.}; no broken links. A leading
\code{\textasciitilde} is the home folder.
\begin{params}
\item[prefix] The text typed; empty: the \code{start} folder.
\end{params}
\returns{\code{\{"items": [\{path, type\}]\}}: up to 50, folders ending with the separator;
none for a folder outside the data roots.}

\apientry{GET api/v1/scan}{What opening a folder would read}
\begin{lstlisting}[style=signature]
GET api/v1/scan?path=PATH
\end{lstlisting}
What opening a folder or a WFDB record would read, from the headers alone, within the limits on
a folder of records: any folder (its records and those of its subfolders, as one session) or a
record (\cref{chap:recordings}). The Open dialog shows it before a folder is opened.
\begin{params}
\item[path] The folder, or the record's \file{.hea} file.
\end{params}
\returns{\code{path}, \code{format}, \code{records}, \code{segments} (of the multi-segment
records), \code{signals} (their names), \code{duration\_s} (the records' total length),
\code{t\_range} (\code{[first, last]} of the records with a stated time, or \code{null}),
\code{untimed} (the number of records without one), \code{bytes} (the size of the records'
files), \code{skipped} (the number of records that cannot be read) and \code{scan\_s} (how long
the scan took).}
\raises{\code{400} neither a folder nor a WFDB record; \code{403}, \code{404} as for
\code{browse}; \code{422} a folder over a limit (more than \code{max\_records} records, records
nested more than \code{max\_depth} levels deep, more than \code{max\_scan\_s} seconds to count
them), or without records.}
\begin{lstlisting}[language=JSON]
{"path": "/data/icu/demo/bed01", "format": "WFDB collection", "records": 24, "segments": 0,
 "signals": ["II", "V", "PLETH", "RESP"], "duration_s": 86400.0,
 "t_range": [1767225600.0, 1767312000.0], "untimed": 0, "bytes": 412318720, "skipped": 0,
 "scan_s": 0.04}
\end{lstlisting}

\apientry{POST api/v1/open}{Open a recording}
\begin{lstlisting}[style=signature]
POST api/v1/open   {"path": PATH, "reindex": BOOL}
\end{lstlisting}
Opens a recording: a WFDB record (its \file{.hea} file), a folder (its records and those of its
subfolders as one session, or a WFDB collection), an EDF or BDF file, or a file an adapter
reads. When its index
exists, the dataset is ready at once; otherwise its index is built in the background, at most
\code{index\_workers} at a time (\cref{chap:settings-reference}). A folder that is not a
collection is first checked against the limits, as \code{scan} does. Opening a recording whose
index is being built answers \code{indexing} again, without starting a second one.
\begin{params}
\item[path] The recording.
\item[reindex] \code{true}: build the index again even if it exists. Default \code{false}.
\end{params}
\returns{\code{\{key, state\}}: \code{state} \code{ready}, or \code{indexing} (follow it with
\code{GET api/v1/datasets/\{key\}}). If building the index fails in any way, the dataset's
state becomes \code{error}, never stays \code{indexing}.}
\raises{\code{400} a body that is not a JSON object, a field of the wrong type, or a path that
is not a recording the viewer reads; \code{403}, \code{404} as for \code{browse}, and
\code{403} for a WFDB record any of whose files (its header, the signal files it names, its
segments, its \file{.mrk} file) leads outside the data roots
(\cref{sec:formats-reference-roots}); \code{413} a body over 64\,KiB; \code{415} a body that is
not \code{application/json}; \code{422} a folder over a limit, as for \code{scan}.}
\begin{lstlisting}[style=shell]
curl -s -H 'Content-Type: application/json' -d '{"path": "/data/icu/demo/rec0001.hea"}' \
  http://localhost:8060/api/v1/open
\end{lstlisting}
\begin{lstlisting}[language=JSON]
{"key": "3fa94c0d2b7e18a65c91", "state": "indexing"}
\end{lstlisting}

\apientry{POST api/v1/cancel}{Stop building an index}
\begin{lstlisting}[style=signature]
POST api/v1/cancel   {"key": KEY}
\end{lstlisting}
Stops the index being built under a key. The dataset's state then becomes \code{cancelled},
and nothing of it is left in the index folder.
\begin{params}
\item[key] The dataset's key, as \code{open} answered it.
\end{params}
\returns{\code{\{key, state: "cancelling"\}}; \code{\{key, state: "cancelled"\}} for an index
whose building was already cancelled.}
\raises{\code{404} no index being built under that key, nor cancelled.}

\apientry{GET api/v1/datasets/\{key\}}{A dataset}
\begin{lstlisting}[style=signature]
GET api/v1/datasets/{key}
\end{lstlisting}
A dataset's state and, once ready, its description.
\returns{one of:
\code{\{state: "indexing", progress: \{state, fraction, elapsed\_s\}\}} (\code{fraction} from
0 to 1); \code{\{state: "error", error\}}; \code{\{state: "cancelled"\}};
\code{\{state: "ready", dataset\}}, where \code{dataset} has its \code{key}, \code{name},
\code{path}, \code{kind}, \code{format} and \code{description}, \code{t\_range}
(\code{[first, last]}), \code{time\_resolution\_s}, \code{channels}, \code{pieces} (each
\code{\{id, ch, t0, fs, n, timed\}}), \code{vitals} (the vital-sign series, each
\code{\{name, unit, n\}}), \code{alarms}, \code{gaps} (the breaks), \code{untimed} (the untimed
spans), \code{late\_blocks} (late data blocks), \code{markers} (\code{\{names, count\}}),
\code{skipped}, \code{notes} and \code{indexed\_s}.}
\raises{\code{404} no dataset indexed under that key; \code{403} a dataset whose recording lies
outside the data roots and was not opened through this server.}

\apientry{GET api/v1/datasets/\{key\}/window}{Waveform samples}
\begin{lstlisting}[style=signature]
GET api/v1/datasets/{key}/window?t0=T&t1=T&px=N&ch=0,1
\end{lstlisting}
The samples of channels over a time span, at about \code{px} points across: the samples
themselves when there are few enough, else the minimum and maximum of each bucket (an
envelope), from the index's levels.
\begin{params}
\item[t0, t1] The time span (required; \code{t1} after \code{t0}).
\item[px] The width in points (default 1200), clamped to 16 to 8000.
\item[ch] The channels, comma-separated channel numbers (none: no channel).
\end{params}
\returns{\code{\{t0, t1, channels\}}: \code{channels} gives, for each channel asked (by its
number, as text), the blocks of its pieces that cover \code{[t0, t1]}, in order of piece id,
each \code{\{piece, kind: "raw", t0, dt, y\}} (samples \code{y}, the first at \code{t0}, then
every \code{dt} seconds) or \code{\{piece, kind: "env", t0, dt, min, max\}} (the minimum and
maximum of each bucket of \code{dt} seconds). When more than \code{px} pieces of a channel fall
in the window, only the first piece (in order of \code{t0}, then id) to start in each pixel
column is answered, so a channel never has more than \code{px} blocks.}
\raises{\code{400} \code{t0} or \code{t1} missing or not numbers, \code{t1} not after
\code{t0}, \code{ch} not a list of the dataset's channel ids (\cref{sec:viewer-api-checks}).}
\begin{lstlisting}[language=JSON]
{"t0": 1767225600.0, "t1": 1767225610.0,
 "channels": {"0": [{"piece": 0, "kind": "raw", "t0": 1767225600.0, "dt": 0.004,
                     "y": [0.12, 0.15, 0.21]}]}}
\end{lstlisting}

\apientry{GET api/v1/datasets/\{key\}/vitals}{Vital-sign series}
\begin{lstlisting}[style=signature]
GET api/v1/datasets/{key}/vitals?t0=T&t1=T&px=N&names=HR|SpO2
\end{lstlisting}
The values of vital-sign series over a time span.
\begin{params}
\item[t0, t1] The time span (required).
\item[px] The width in pixels (default 1200), clamped to 16 to 8000.
\item[names] The series, separated by \code{\textbar}; a name the dataset does not have is left
  out.
\end{params}
\returns{an object with one entry per series: \code{\{unit, kind: "points", t, v\}} (its
points in the span), or, when there are more than three points per pixel column,
\code{\{unit, kind: "env", t, min, max\}} (the minimum and maximum of each column, at the
column's first time).}

\apientry{GET api/v1/datasets/\{key\}/vitals\_at}{The vital signs at a time}
\begin{lstlisting}[style=signature]
GET api/v1/datasets/{key}/vitals_at?t=T
\end{lstlisting}
The latest value of every vital sign at a time, no older than 10 minutes.
\begin{params}
\item[t] The time (required).
\end{params}
\returns{a list of \code{\{name, unit, value, t\}}.}
\begin{lstlisting}[language=JSON]
[{"name": "HR", "unit": "bpm", "value": 82.0, "t": 1767225598.0},
 {"name": "SpO2", "unit": "%", "value": 97.0, "t": 1767225590.0}]
\end{lstlisting}

\apientry{GET api/v1/datasets/\{key\}/markers}{Waveform markers}
\begin{lstlisting}[style=signature]
GET api/v1/datasets/{key}/markers?t0=T&t1=T&ch=0,1
\end{lstlisting}
The waveform markers (beat labels, events marked on a signal) of channels over a time span.
\begin{params}
\item[t0, t1] The time span (required).
\item[ch] The channels, comma-separated.
\end{params}
\returns{\code{\{names, items, truncated\}}: each item \code{[channel, time, name index]}, the
index into \code{names}; at most 20\,000 items, and then \code{truncated} is \code{true}.}

\apientry{GET api/v1/docs}{Interactive documentation}
\begin{lstlisting}[style=signature]
GET api/v1/docs
GET api/v1/openapi.json
\end{lstlisting}
The interactive documentation of the API, and its OpenAPI description: the Python viewer only.

\subsection{Through Medorama}
\label{sec:viewer-api-gateway}

\md{} runs a VitalStream viewer under its own address, \file{/viewer/}, behind its sign-in
(\cref{chap:viewers-in-medorama}). Its gateway passes on only the requests that the page needs,
each with its rule (\cref{tab:viewer-api-gateway}); it refuses everything else (\code{404}
\code{not available through the catalog}), the viewer's own \code{login} and \code{logout}
included. Each request passed on carries the viewer's sign-in cookie, made by \md{}; the
user's \md{} cookie is never passed on. A change of the viewer API must be given a rule in the
gateway, or \md{} refuses the new request.

\begin{xltabular}{\linewidth}{@{}P{4.2cm}L@{}}
\caption{The requests that \md{}'s gateway passes to VitalStream, and their rules.}
\label{tab:viewer-api-gateway}\\
\toprule
{\tabhead Request} & {\tabhead Rule} \\
\midrule
\endfirsthead
\toprule
{\tabhead Request} & {\tabhead Rule} \\
\midrule
\endhead
\bottomrule
\endfoot
\code{GET /}, \code{GET static/...} & The page, for signed-in users; without a session, a
  redirection to \md{}'s page. \\
\code{GET api/v1/config} & Passes. For users who are not admins, \code{places}, \code{start}
  and \code{home} point into the folders of the catalog's recordings, and \code{data\_roots} is
  empty; in Privacy Mode all four are empty. \\
\code{GET api/v1/browse}, \code{GET api/v1/complete} & Admins: inside \md{}'s
  \code{data\_roots} (anywhere when it is empty). Other users: only the folders that lead to
  recordings in the catalog, showing only those recordings (\code{browse} needs a
  \code{path}). Refused in Privacy Mode (\code{403}: folder and file names may identify). \\
\code{GET api/v1/scan} & As \code{open}; refused in Privacy Mode. \\
\code{POST api/v1/open} & Admins: inside \code{data\_roots}; other users: only recordings in
  the catalog. In Privacy Mode, a recording is opened by its token (\code{mo-rec:<id>}), never
  by a path. Audited (\code{view}). A JSON body is required (\code{415} otherwise). \\
\code{POST api/v1/cancel} & Only the indexes of datasets the user opened through the gateway. \\
\code{GET api/v1/datasets/\{key\}...} & Only the datasets the user opened through the gateway.
  In Privacy Mode, the times the page sends (\code{t0}, \code{t1}, \code{t}) are moved forward
  by the recording's offset before they reach the viewer, the dataset's \code{path} and
  \code{name} are replaced (the token, \code{[Recording]}), and the times of \code{t\_range}
  and of the markers moved back. The texts of the dataset's annotations (its events of the kind
  \code{annotation}, from EDF+ files) are treated like the text of notes: they read
  \code{[Annotation]} when Privacy Mode withholds notes' text, and otherwise have the patient's
  names and date of birth replaced (\code{gate\_annotations} in \file{medorama/src/gateway.rs};
  \cref{sec:protecting-vitalstream}). \\
\end{xltabular}

Every answer then passes the privacy gate like any other answer of \md{}
(\cref{chap:protecting}): in Privacy Mode the times in its JSON move back and identifiers are
replaced, the viewer's own error texts (which name files) are replaced by a fixed sentence, and
the records of a folder it left out (\code{skipped}) are named by number. A request body larger
than 1\,MiB is refused (\code{413}). A request through the gateway that changes something
(\code{open}, \code{cancel}) must carry \code{X-Medorama: 1}, as every changing request to
\md{} does (the viewer's page sends it), and be JSON (\code{Content-Type: application/json}),
which another site's page cannot send unasked. The viewer's page under \file{/viewer/} is served
with a Content Security Policy of its own scripts and styles only (\code{default-src 'self';
script-src 'self'; img-src 'self' data:; object-src 'none'; base-uri 'none'; frame-ancestors
'self'}).

\section{The addresses of the viewers}
\label{sec:viewer-api-addresses}

A viewer is told what to show by the fragment of its address (after \code{\#}), so that another
program can open it at a file (\cref{tab:viewer-api-fragments}). Changing the fragment shows
another file without reloading the page. The files are fetched with the page's own credentials
(same origin), so that a viewer served by \md{} reads exactly what \md{} lets its user read,
through \code{api/v1/patients/\{id\}/file/\{item\}/...} (\cref{chap:http-api}).

\begin{xltabular}{\linewidth}{@{}P{1.9cm}P{4.0cm}L@{}}
\caption{The fragments of the viewers' addresses.}
\label{tab:viewer-api-fragments}\\
\toprule
{\tabhead Viewer} & {\tabhead Fragment} & {\tabhead Means} \\
\midrule
\endfirsthead
\toprule
{\tabhead Viewer} & {\tabhead Fragment} & {\tabhead Means} \\
\midrule
\endhead
\bottomrule
\endfoot
ECG & \code{src=URL} & The file: a WFDB header (its signal files, and a multi-segment record's
  segments, are read beside it: \code{URL}'s folder followed by their names), a MUSE XML or
  DICOM ECG, a JSON ECG, a picture or a PDF. \\
PACS & \code{src=URL} & One file. \\
PACS & \code{list=URL} & A folder: \code{URL} answers \code{\{"files": ["IM0001.dcm", ...]\}},
  and each file is read at \code{URL} followed by its name. \\
PACS & \code{mm=...}, \code{fps=...} & The millimeters per pixel of pictures and videos; the
  frame rate of pictures (one loop). \\
ECG & \code{shift=...}, \code{cap90=1} & For a host that de-identifies in the page: the
  report's dates and the record's start moved back, ages above 89 years shown as 90+ (\md{}
  de-identifies on its server
  instead). \\
ECG, PACS & \code{name=...} & The name to show. \\
All three & \code{theme=light} or \code{dark} & The colors, as the host's: applied, not
  remembered as the viewer's own choice (VitalStream's theme button still switches them).
  \md{} passes its page's theme to every viewer tab, when it opens and when it changes. \\
All three & \code{who=...}, \code{what=...}, \code{pos=...}, \code{nav=back,prev,next} & A host's
  context: whose file it is and what it is, shown in place of the file's name, with the buttons
  \ui{Patient} and \ui{\textless}~\ui{\textgreater}, which the viewer only reports to the host
  page (\code{postMessage}, same origin). Without them, the file's name shows. \\
ECG & \code{compact=1} & The paper alone, fitted to the frame (shown inside another page). \\
VitalStream & \code{path=...\&t0=...\&t1=...} & A recording on the server (its path), at a time
  span (Unix seconds). \\
VitalStream & \code{browse=FOLDER} & The file browser, open at a folder. \\
\end{xltabular}

For example, \md{}'s patient page draws an ECG on paper with the ECG viewer at
\code{viewers/ecg/\#src=/api/v1/patients/3/file/52\&name=ecg.xml\&theme=light\&compact=1}
(the values URL-encoded), and VitalStream opens a recording at a span with
\code{\#path=/data/icu/demo/rec0001.hea\&t0=1767225600\&t1=1767229200} (inside \md{}, with the
host's context and \code{theme=} after it).

\section{The ECG and PACS viewers' server}
\label{sec:viewer-api-serve}

The ECG and PACS viewers are pages that read in the browser the files they are given. The
command \cmd{serve} of their Python packages, \cmd{muse-py serve} (the ECG viewer, port 8080 by
default) and \cmd{dicom-py serve} (the PACS viewer, port 8081 by default), also serves the
page from a computer together with the files under a folder of it, the \emph{root}: on a
server where the data are, reached through an SSH tunnel, as for VitalStream. Its options are
in \cref{chap:cli}. The server uses only the Python standard library; it is the same in both
packages.

\begin{itemize}
\item \textbf{Root.} Only the files and folders under the root are listed and served, every
  symbolic link resolved (a link that leads outside is refused); hidden entries (names that
  start with \code{.}) are neither listed nor served. The root is \opt{-{}-root}
  (default: the home folder), and lies inside the root the package was installed with
  (\code{VIEWER\_ROOT} at installation; default: the whole file system), never outside it.
\item \textbf{Host.} A request must name the server by a loopback name, an IP address or this
  host (\opt{-{}-host}); any other name is refused (\code{421}; no DNS rebinding).
\item \textbf{Password} (optional): the same file and scheme as VitalStream's, so one password
  serves all the viewers. It is set once with \opt{-{}-set-password} (at least 8 characters);
  the page then asks for it, in the cookie \code{viewer\_<viewer>\_<port>}
  (\code{viewer\_ecg\_8080}). For \cmd{muse-py serve} and \cmd{dicom-py serve}, a sign-in lasts
  30 days, until the server stops, or until a browser signs out, which ends every sign-in; after
  a wrong password the next try, from any browser, waits 1\,s, then 2, 4 \dots{} up to 60\,s.
  The cookie is signed with a key of the running server too: browsers send a cookie to every
  port of a host, and another server on 127.0.0.1 cannot use it. A package installed with
  \code{VIEWER\_PASSWORD=0} has no password support: it ignores a password already set in the
  file, never asks for one, and \opt{-{}-set-password} and \opt{-{}-clear-password} fail
  (\code{this viewer was installed without password support}). Without a password, anyone who
  can connect to the port can read the files under the root; the server says so when it starts.
\item \textbf{Requests} (\cref{tab:viewer-api-serve}). Without the sign-in cookie (a password
  being set), \code{GET /} answers the sign-in form and every other request \code{401}
  \code{\{"detail": "sign in required"\}}. Errors are \code{\{"detail": "..."\}}. The server
  keeps no log of the paths asked for.
\item \textbf{Files} are sent as they are read, in part when a browser asks for a range
  (\code{Range}), under a policy that runs nothing in them (\code{Content-Security-Policy:
  sandbox}); a kind of file the viewers do not read is sent as a download.
\end{itemize}

\begin{xltabular}{\linewidth}{@{}P{3.6cm}L@{}}
\caption{The requests of \cmd{muse-py serve} and \cmd{dicom-py serve}.}
\label{tab:viewer-api-serve}\\
\toprule
{\tabhead Request} & {\tabhead Answer} \\
\midrule
\endfirsthead
\toprule
{\tabhead Request} & {\tabhead Answer} \\
\midrule
\endhead
\bottomrule
\endfoot
\code{GET /} & The viewer's page (or, before sign-in, the sign-in form, \code{401}); its other
  files by their names (\code{GET app.js}). \\
\code{GET browse?path=REL} & A folder under the root, as an HTML page (\ui{server folders}):
  its subfolders and the files the viewer opens, each a link that opens it in the page
  (\code{./\#src=files/REL\&name=NAME}). The ECG viewer lists \file{.hea}, \file{.xml},
  \file{.json}, \file{.dcm}, \file{.pdf}, \file{.png}, \file{.jpg} and \file{.jpeg} files; the
  PACS viewer \file{.dcm}, \file{.dicom}, \file{.mp4}, \file{.webm}, \file{.mov}, \file{.png},
  \file{.jpg}, \file{.jpeg}, \file{.pdf} and files without an extension, with a link
  \ui{Open the records in this folder} (\code{\#list=list/REL/}) when the folder has such files.
  Hidden names are left out; at most 5000 entries, and the page says so when there are more.
  \code{403} for a path that is not a folder
  under the root. A file or a folder opened from this list is shown with \code{\#src=} (or
  \code{\#list=}), where the page hides its \ui{Server folders\dots} button: the browser's Back
  button returns to the list. \\
\code{GET files/REL} & A file under the root, whole or the range asked (\code{Range}).
  \code{403} outside the root or hidden, \code{404} not a file. \\
\code{GET list/REL/} & A folder's files (not its subfolders, not hidden ones, at most 5000) as
  \code{\{"files": [NAME...]\}}, with \code{"truncated": true} when there are more, for the PACS
  viewer's \code{\#list=}. Without the final
  \code{/}, the file \code{REL}, as \code{files/}. \\
\code{GET served} & \code{\{"viewer": "ecg" | "pacs", "root": FOLDER\}}: tells the page that it
  is served this way, so that it shows its \ui{Server folders\dots} button. \\
\code{POST login} & The sign-in form's password (form field \code{password}, at most 4\,KiB):
  sets the cookie and redirects (\code{303}) to \code{./}; a wrong password answers the form
  again with \code{401} and \code{Retry-After}, the wait before the next try (1\,s, then 2,
  4 \dots{} up to 60\,s); a password sent before the wait has passed is not checked
  (\code{429}). \\
\code{POST logout} & Removes the cookie, ends every sign-in, and redirects to \code{./}. \\
\end{xltabular}

\code{REL} is a path relative to the root, with \code{/} between folders.
\begin{lstlisting}[style=shell]
muse-py serve --root /data/ecg --no-browser
curl -s http://127.0.0.1:8080/served
dicom-py serve --root /data/pacs --port 8081
curl -s http://127.0.0.1:8081/list/study01/
\end{lstlisting}
